% =============================================================================
% EXPERIMENTS - CITA Paper in ALKALI Format
% =============================================================================



\vspace{-2mm}
\subsection{Training Pipeline}
\label{sec:training_pipeline_exp}

Following the staged approach from Section~\ref{sec:methodology}, we train on PKU-SafeRLHF~\cite{ji2024pku,bai2022training} using NVIDIA A100 GPUs with LoRA~\cite{hu2022lora} adapters (SFT/DPO/CITA on A100-40GB; PPO/GRPO on A100-80GB). Training proceeds through four stages:


\begin{figure}[t]
\centering
\includegraphics[width=\columnwidth]{figures/pipeline/training_pipeline.pdf}
\caption{\textbf{Training Pipeline}: All methods branch from SFT. PPO and GRPO are online methods requiring reward model/functions. DPO uses offline preference pairs. CITA stacks on DPO with instruction-conditioning and mandatory KL regularization.}
\label{fig:training_pipeline}
\vspace{-1em}
\end{figure}

\begin{enumerate}[leftmargin=*,itemsep=0pt]
    \item \textbf{Base}: Llama-3.1-8B~\cite{dubey2024llama} (pretrained weights)
    \item \textbf{SFT}: Fine-tune on PKU chosen responses
    \item \textbf{DPO}: Train on PKU preference pairs
    \item \textbf{CITA}: Add instruction-conditioning with mandatory KL
\end{enumerate}



\vspace{-2mm}
\subsection*{Model Variants}
\label{sec:model_variants}

We train 10 model variants across 5 methods (SFT, PPO, GRPO, DPO, CITA) $\times$ 2 instruction settings (NoInstruct, Instruct).

% \begin{table}[H]
% \centering
% \footnotesize
% \begin{tabular}{@{}lll@{}}
% \toprule
% \textbf{Model} & \textbf{Training Path} & \textbf{Instruction} \\
% \midrule
% SFT\_NI & Base$\rightarrow$SFT & No \\
% SFT\_I & Base$\rightarrow$SFT & Yes \\
% \midrule
% PPO\_NI & Base$\rightarrow$SFT$\rightarrow$PPO & No \\
% PPO\_I & Base$\rightarrow$SFT$\rightarrow$PPO & Yes \\
% \midrule
% GRPO\_NI & Base$\rightarrow$SFT$\rightarrow$GRPO & No \\
% GRPO\_I & Base$\rightarrow$SFT$\rightarrow$GRPO & Yes \\
% \midrule
% DPO\_NI & Base$\rightarrow$SFT$\rightarrow$DPO & No \\
% DPO\_I & Base$\rightarrow$SFT$\rightarrow$DPO & Yes \\
% \midrule
% CITA\_NI & Base$\rightarrow$SFT$\rightarrow$DPO$\rightarrow$CITA & No \\
% CITA\_I & Base$\rightarrow$SFT$\rightarrow$DPO$\rightarrow$CITA & Yes \\
% \bottomrule
% \end{tabular}
% \caption{10 model variants. NI = NoInstruct (standard prompts), I = Instruct (with behavioral instructions). PPO, GRPO, DPO branch from SFT; CITA stacks on DPO.}
% \label{tab:model_variants}
% \vspace{-4mm}
% \end{table}

\textbf{Comparison Design.} By evaluating NoInstruct vs. Instruct variants, we measure each method's \textit{instruction sensitivity}---the improvement from adding alignment instructions.

\vspace{-2mm}
\subsection*{Hyperparameter Optimization}
\label{sec:hyperparams}

We use Optuna~\cite{akiba2019optuna} with TPE sampler (12 trials) to optimize hyperparameters, following best practices from hyperparameter optimization literature~\cite{liang2023holistic}. Key finding: \textbf{CITA\_Instruct requires $\sim$50\% lower learning rate} than CITA\_NoInstruct because instruction-augmented sequences are 30--40\% longer, causing larger gradients.

\begin{table}[H]
\centering
\footnotesize
\setlength{\tabcolsep}{6pt}
\renewcommand{\arraystretch}{1.2}
\begin{tabular}{@{}lcc@{}}
\toprule
\textbf{Hyperparameter} & \textbf{NoInstruct (T5)} & \textbf{Instruct (T7)} \\
\midrule
Learning rate & 6.83e-6 & 5.41e-6 \\ \hline
$\lambda_{\text{KL}}$ & 0.00052 & 0.00023 \\ \hline
$\beta$ (DPO temp.) & 0.119 & 0.107 \\ \hline
Weight decay & 0.0091 & 0.0109 \\ \hline
Warmup ratio & 7.5\% & 10.0\% \\
\midrule
Final margin & 6.95 & 7.52 \\ \hline
Accuracy & 89.5\% & 89.0\% \\
\bottomrule
\end{tabular}
\caption{Optimal hyperparameters from Optuna search. T5/T7 = trial number of best configuration.}
\label{tab:hyperparams_exp}
\vspace{-4mm}
\end{table}

\vspace{-2mm}
\subsection*{Training epochs}
\label{sec:training_dynamics}

We analyze key training curves across all alignment methods. Training accuracy metrics (reward accuracy for DPO/CITA, mean token accuracy for SFT) are provided in Appendix~\ref{sec:appendix_training_curves}, Figure~\ref{fig:training_accuracy_appendix}.



% Training accuracy figure moved to Appendix (see Figure~\ref{fig:training_accuracy_appendix})
% for detailed 1x2 subplot comparing DPO/CITA reward accuracy and SFT mean token accuracy

\begin{figure}[H]
\centering
\includegraphics[width=0.65\columnwidth]{figures/training/dpo_cita_margins.pdf}
\caption{\textbf{Reward Margins}: CITA\_Instruct (7.5) $>$ CITA\_NoInstruct (7.2) $>$ DPO ($\sim$6.0). Higher margins indicate stronger preference learning.}
\label{fig:training_margins}
\vspace{-3mm}
\end{figure}

\textbf{Key Observation.} CITA achieves higher reward margins (7.5 vs. 6.0) than DPO despite similar accuracy, validating that mandatory KL regularization leads to stronger preference differentiation.

